\textbf{Phần 1. Trắc nghiệm lựa chọn}
\setcounterpageref{ex}{0}

\begin{ex} %[4D10N54]
Thương của phép chia $3 : 4$ được viết dưới dạng phân số là
\choice
{$\dfrac{4}{3}$}
{$3{,}4$}
{$4{,}3$}
{\True $\dfrac{3}{4}$}
\loigiai{
$3 : 4 = \dfrac{3}{4}$.\\
\textbf{Đáp án của câu hỏi là D}
}
\end{ex}

\begin{ex} %[4D10N54]
Phân số $\dfrac{5}{1}$ có giá trị bằng
\choice
{$1$}
{\True $5$}
{$0$}
{$10$}
\loigiai{
$\dfrac{5}{1} = 5 : 1 = 5$.\\
\textbf{Đáp án của câu hỏi là B}
}
\end{ex}

\begin{ex} %[4D10H54]
Phân số $\dfrac{0}{8}$ có giá trị bằng
\choice
{$8$}
{$1$}
{\True $0$}
{Không xác định}
\loigiai{
$\dfrac{0}{8} = 0 : 8 = 0$.\\
\textbf{Đáp án của câu hỏi là C}
}
\end{ex}

\begin{ex} %[4D10H54]
Mỗi số tự nhiên $a$ đều có thể viết dưới dạng phân số có mẫu số bằng 1 là
\choice
{\True $\dfrac{a}{1}$}
{$\dfrac{1}{a}$}
{$\dfrac{a}{a}$}
{$\dfrac{0}{a}$}
\loigiai{
$a = \dfrac{a}{1}$.\\
\textbf{Đáp án của câu hỏi là A}
}
\end{ex}

\begin{ex} %[4D10H54]
Chia đều 5 cái bánh cho 4 bạn. Mỗi bạn nhận được số phần cái bánh là
\choice
{$\dfrac{4}{5}$ cái bánh}
{$1$ cái bánh}
{$2$ cái bánh}
{\True $\dfrac{5}{4}$ cái bánh}
\loigiai{
$5 : 4 = \dfrac{5}{4}$ cái bánh.\\
\textbf{Đáp án của câu hỏi là D}
}
\end{ex}

\begin{ex} %[4D10H54]
Phân số nào sau đây lớn hơn 1?
\choice
{$\dfrac{3}{5}$}
{$\dfrac{4}{4}$}
{\True $\dfrac{7}{4}$}
{$\dfrac{1}{2}$}
\loigiai{
Phân số có tử số lớn hơn mẫu số thì lớn hơn 1: $\dfrac{7}{4} > 1$.\\
\textbf{Đáp án của câu hỏi là C}
}
\end{ex}

\begin{ex} %[4D10V54]
Phân số nào sau đây bằng 1?
\choice
{\True $\dfrac{8}{8}$}
{$\dfrac{5}{6}$}
{$\dfrac{9}{7}$}
{$\dfrac{0}{8}$}
\loigiai{
Tử số bằng mẫu số thì phân số bằng 1: $\dfrac{8}{8} = 1$.\\
\textbf{Đáp án của câu hỏi là A}
}
\end{ex}

\begin{ex} %[4D10V54]
Có 3 lít nước rót đều vào 5 chai. Mỗi chai có số lít nước là
\choice
{$\dfrac{5}{3}\text{ l}$}
{$1\text{ l}$}
{\True $\dfrac{3}{5}\text{ l}$}
{$0{,}5\text{ l}$}
\loigiai{
$3 : 5 = \dfrac{3}{5}\text{ l}$.\\
\textbf{Đáp án của câu hỏi là C}
}
\end{ex}

\begin{ex} %[4D10V54]
Trong các phân số sau, phân số nào bé hơn 1?
\choice
{$\dfrac{4}{3}$}
{\True $\dfrac{5}{7}$}
{$\dfrac{9}{9}$}
{$\dfrac{11}{8}$}
\loigiai{
Tử số nhỏ hơn mẫu số: $\dfrac{5}{7} < 1$.\\
\textbf{Đáp án của câu hỏi là B}
}
\end{ex}

\begin{ex} %[4D10V54]
Viết thương $7 : 10$ dưới dạng phân số:
\choice
{$\dfrac{10}{7}$}
{$0{,}7$}
{$70$}
{\True $\dfrac{7}{10}$}
\loigiai{
$7 : 10 = \dfrac{7}{10}$.\\
\textbf{Đáp án của câu hỏi là D}
}
\end{ex}

\begin{ex} %[4D10V54]
Chia đều $12\text{ kg}$ gạo vào 5 túi. Mỗi túi có bao nhiêu ki-lô-gam gạo?
\choice
{$\dfrac{5}{12}\text{ kg}$}
{\True $\dfrac{12}{5}\text{ kg}$ ($2{,}4\text{ kg}$)}
{$2\text{ kg}$}
{$3\text{ kg}$}
\loigiai{
$12 : 5 = \dfrac{12}{5}\text{ kg}$.\\
\textbf{Đáp án của câu hỏi là B}
}
\end{ex}

\begin{ex} %[4D10C54]
Có bao nhiêu phân số có tử số và mẫu số là các số tự nhiên có một chữ số và có giá trị bằng 1?
\choice
{\True $9$ phân số}
{$8$ phân số}
{$10$ phân số}
{$7$ phân số}
\loigiai{
Các phân số là: $\dfrac{1}{1}, \dfrac{2}{2}, \dfrac{3}{3}, \dfrac{4}{4}, \dfrac{5}{5}, \dfrac{6}{6}, \dfrac{7}{7}, \dfrac{8}{8}, \dfrac{9}{9}$ (9 phân số, mẫu số khác 0).\\
\textbf{Đáp án của câu hỏi là A}
}
\end{ex}

\textbf{Phần 2. Trắc nghiệm đúng sai}
\setcounterpageref{ex}{0}

\begin{ex} %[4D10H54]
Xét tính đúng sai về phân số và phép chia số tự nhiên:
\choiceTF
{\True Thương của phép chia số tự nhiên cho số tự nhiên khác 0 có thể viết dưới dạng phân số.}
{\True Phân số có tử số lớn hơn mẫu số thì phân số đó lớn hơn 1.}
{Phân số $\dfrac{6}{6}$ bé hơn 1.}
{\True Số 0 có thể viết dưới dạng phân số là $\dfrac{0}{5}$ (hoặc mẫu số bất kì khác 0).}
\loigiai{
\begin{itemize}
    \item Thương viết thành phân số: Đúng.
    \item Tử $>$ mẫu $\Rightarrow > 1$: Đúng.
    \item $\dfrac{6}{6} = 1$ (không bé hơn 1). Sai.
    \item $0 = \dfrac{0}{5}$: Đúng.
\end{itemize}
\begin{itemize}
    \item a) \textbf{Đúng}.
    \item b) \textbf{Đúng}.
    \item c) \textbf{Sai}.
    \item d) \textbf{Đúng}.
\end{itemize}
}
\end{ex}

\begin{ex} %[4D10V54]
Cho các phân số: $\dfrac{3}{4}, \dfrac{5}{5}, \dfrac{7}{2}, \dfrac{9}{10}, \dfrac{12}{11}$. Xét tính đúng sai:
\choiceTF
{\True Có 2 phân số bé hơn 1 (là $\dfrac{3}{4}$ và $\dfrac{9}{10}$).}
{\True Có 1 phân số bằng 1 (là $\dfrac{5}{5}$).}
{Có 3 phân số lớn hơn 1.}
{\True Có 2 phân số lớn hơn 1 (là $\dfrac{7}{2}$ và $\dfrac{12}{11}$).}
\loigiai{
\begin{itemize}
    \item Bé hơn 1: $\dfrac{3}{4}, \dfrac{9}{10}$ (2 phân số). Đúng.
    \item Bằng 1: $\dfrac{5}{5}$ (1 phân số). Đúng.
    \item Lớn hơn 1: $\dfrac{7}{2}, \dfrac{12}{11}$ (chỉ có 2 phân số, không phải 3). Ý 3 Sai, Ý 4 Đúng.
\end{itemize}
\begin{itemize}
    \item a) \textbf{Đúng}.
    \item b) \textbf{Đúng}.
    \item c) \textbf{Sai}.
    \item d) \textbf{Đúng}.
\end{itemize}
}
\end{ex}

\textbf{Phần 3. Trả lời ngắn}
\setcounterpageref{ex}{0}

\begin{ex} %[4D10H54]
Viết thương $4 : 7$ dưới dạng phân số (dạng a/b).
\shortans{4/7}
\loigiai{
\textbf{Đáp số:} $4/7$\\
$4/7$.
}
\end{ex}

\begin{ex} %[4D10H54]
Viết số tự nhiên 9 dưới dạng phân số có mẫu số bằng 1 (dạng a/b).
\shortans{9/1}
\loigiai{
\textbf{Đáp số:} $9/1$\\
$9/1$.
}
\end{ex}

\begin{ex} %[4D10V54]
Chia 7 quả cam cho 4 người. Viết số cam mỗi người nhận được dưới dạng phân số (dạng a/b).
\shortans{7/4}
\loigiai{
\textbf{Đáp số:} $7/4$\\
$7/4$ quả cam.
}
\end{ex}

\begin{ex} %[4D10V54]
Có bao nhiêu phân số nhỏ hơn 1 có mẫu số bằng 5 (tử số là số tự nhiên khác 0)?
\shortans{4}
\loigiai{
\textbf{Đáp số:} $4$\\
$1/5, 2/5, 3/5, 4/5$ (4 phân số).
}
\end{ex}

\begin{ex} %[4D10V54]
Tìm số tự nhiên $x$ biết $\dfrac{x}{8} = 1$.
\shortans{8}
\loigiai{
\textbf{Đáp số:} $8$\\
$x = 8$.
}
\end{ex}

\begin{ex} %[4D10C54]
Có 5 đoạn dây, mỗi đoạn dài $2\text{ m}$. Người ta nối lại rồi chia đều thành 8 phần bằng nhau. Viết độ dài mỗi phần theo mét dưới dạng phân số rút gọn (dạng a/b).
\shortans{5/4}
\loigiai{
\textbf{Đáp số:} $5/4$\\
Tổng chiều dài: $5 \times 2 = 10\text{ m}$. Mỗi phần: $10 : 8 = 10/8 = 5/4\text{ m}$.
}
\end{ex}
